% 4.6 由小群诱导的表示的实性
\section{由小群诱导的表示的实性}

参照定义 1.3.7 与定理 1.3.9，$\mathbf{G}$ 的表示按和
\begin{equation*}
\frac{1}{|\mathbf{G}|}\sum^{|\mathbf{G}|}_{j=1}\chi(G^{2}_{j}) =
\left\{
\begin{array}{c}
1\\
-1\\
0
\end{array}
\right.
\tag{4.6.1}
\end{equation*}
的取值分为三种。

现在把这个判据用于表示 $\Delta^{i}_{j} = \Gamma^{i}_{j} \uparrow \mathbf{G}$——它是在 $\mathbf{G}$ 中由小群 $\mathbf{K}^{i}$ 的小表示 $\Gamma^{i}_{j}$ 诱导的。记号与 4.3 节相同。$\mathbf{K}^{i}$ 的元素形如 $r_{\gamma}t_{a}$。

回忆 $\mathbf{G}$ 关于 $\mathbf{K}^{i}$ 的左陪集分解（式 (4.1.1)）：
\begin{equation*}
\mathbf{G} = \sum_{\tau}p_{\tau}\mathbf{K}^{i}.
\tag{4.6.2}
\end{equation*}
还回忆式 (4.1.9) 的内容：
\begin{equation*}
\chi_{\tau}(g^{2}) =
\left\{
\begin{array}{ll}
\chi^{i}_{j}(r_{\gamma}t_{a}) & \text{若 } g^{2} = p_{\tau}r_{\gamma}t_{a}p^{-1}_{\tau} \in \mathbf{K}^{i}_{\tau},\\
0 & \text{否则},
\end{array}
\right.
\tag{4.6.3}
\end{equation*}
以及
\begin{equation*}
\chi(g^{2}) = \sum_{\tau}\chi_{\tau}(g^{2}).
\tag{4.6.4}
\end{equation*}
于是和 (4.6.1) 成为
\begin{equation*}
\frac{1}{|\mathbf{G}|}\sum_{g \in \mathbf{G}}\sum_{\tau}\chi_{\tau}(g^{2}) = \frac{1}{|\mathbf{G}|}\sum_{\tau}\sum_{g \in \mathbf{G}}\chi^{i}_{j}(r_{\gamma}t_{a}),
\tag{4.6.5}
\end{equation*}
其中对 $g$ 的求和现在限制在满足 $g^{2} = p_{\tau}r_{\gamma}t_{a}p^{-1}_{\tau} \in \mathbf{K}^{i}_{\tau}$ 的元素上。现在若 $h^{2} \in \mathbf{K}^{i}$，即 $h^{2} = r_{\gamma}t_{a}$（对某个 $\gamma, a$），则 $g = p_{\tau}h$ 的平方属于 $\mathbf{K}^{i}_{\tau}$，于是式 (4.6.5) 成为
\begin{equation*}
\frac{q_{i}}{|\mathbf{G}|}\sum_{h \in \mathbf{G}}\chi^{i}_{j}(h^{2})
\tag{4.6.6}
\end{equation*}
其中对 $h$ 的求和限制在满足 $h^{2} \in \mathbf{K}^{i}$ 的元素上。

现在若 $h = r_{\alpha}t_{b}$，则 $h^{2} = r_{\alpha\alpha}t_{\alpha\alpha}[t_{b}]_{\alpha}t_{b}$，故 $h^{2} \in \mathbf{K}^{i}$ 当 $\alpha^{2} \in \overline{\mathbf{K}}^{i}$。于是和 (4.6.6) 成为
\begin{equation*}
\frac{q_{i}}{|\mathbf{G}|}\sum_{\alpha^{2} \in \overline{\mathbf{K}}^{i}}\sum_{b}\chi^{i}_{j}(r_{\alpha\alpha}t_{\alpha\alpha}[t_{b}]_{\alpha}t_{b}).
\tag{4.6.7}
\end{equation*}
回忆在 $\Gamma^{i}_{j}$ 中元素 $t$ 的矩阵是标量矩阵 $D^{i}(t)\mathbf{1}$，故和 (4.6.7) 约化为
\begin{equation*}
\frac{q_{i}}{|\mathbf{G}|}\sum_{\alpha^{2} \in \overline{\mathbf{K}}^{i}}\sum_{b}\chi^{i}_{j}(r_{\alpha\alpha})\mathbf{D}^{i}(t_{\alpha\alpha})\mathbf{D}^{i}([t_{b}]_{\alpha})\mathbf{D}^{i}(t_{b}).
\tag{4.6.8}
\end{equation*}
由式 (4.2.17)，$\mathbf{D}^{i}([t_{b}]_{\alpha}) = \mathbf{D}^{i}_{\alpha}(t_{b})$，而和对 $b$ 的 $\sum_{b}\mathbf{D}^{i}_{\alpha}(t_{b})\mathbf{D}^{i}(t_{b})$ 由正交关系将除非 $D^{i}_{\alpha} = (D^{i})^{-1}$ 时为零。于是式 (4.6.8) 成为
\begin{equation*}
\frac{q_{i}}{h}\sum_{\alpha}\chi^{i}_{j}(r^{2}_{\alpha})
\tag{4.6.9}
\end{equation*}
其中式 (4.6.9) 中对 $\alpha$ 的求和限制在那些满足 $D^{i}_{\alpha} = (D^{i})^{-1}$ 的元素 $\alpha \in \mathbf{R}$（$\mathbf{R}$ 是 $\mathbf{G}$ 关于 $\mathbf{T}$ 的陪集代表的下标构成的群；见定理 4.2.2）上。$\Delta^{i}_{j}$ 的种类划分取决于这个和是 $+1$、$-1$ 还是 $0$。可以把这一结果总结为：

\noindent\textbf{定理 4.6.1}\quad 在 $\mathbf{G}$ 中由小群 $\mathbf{K}^{i}$ 的小表示 $\Gamma^{i}_{j}$ 诱导的表示 $\Delta^{i}_{j} = \Gamma^{i}_{j} \uparrow \mathbf{G}$ 属于第一、第二或第三种，按照
\begin{equation*}
\frac{q_{i}}{h}\sum_{\alpha}\chi^{i}_{j}(r^{2}_{\alpha}) = 1,\ -1,\ \text{或}\ 0
\tag{4.6.10}
\end{equation*}
而定，其中对 $\alpha$ 的求和限制在满足 $D^{i}_{\alpha} = (D^{i})^{-1}$ 的那些 $\alpha$ 上。

使用诱导表示的好处是：式 (4.6.10) 涉及的求和元素的个数比直接使用式 (4.6.1) 要少得多。在 2.4 节我们已指出，式 (1.3.12) 或 (4.6.1) 可以直接用于晶体学点群的表示。举一个例子，考虑 $422\ (D_{4})$ 的表示 $E$（见 4.5 节）：
\begin{equation*}
\begin{array}{l}
422 = E.4 + C_{2x}.4\\
\left(D_{4} = E.C_{4} + C_{2x}.C_{4}\right)
\end{array}
\tag{4.5.1}
\end{equation*}
而 $E$ 可以看作由 $4\ (C_{4})$ 的表示 ${}^{1}E$ 在 $422\ (D_{4})$ 中诱导的表示，即 ${}^{1}E \uparrow 422$；因此 $h = 2$。$E$ 的星有两个成员 ${}^{1}E$ 与 ${}^{2}E$（见表 4.2 及 4.5 节正文），故 $q_{i} = 2$。式 (4.6.10) 中的和只需对满足 $D^{i}_{\alpha} = (D^{i})^{-1}$ 的 $r_{\alpha}$ 计算。列出 ${}^{1}E$ 及其共轭表示的数值：
\begin{center}
\begin{tabular}{lccccc}
 & & $E$ & $C^{+}_{4z}$ & $C_{2z}$ & $C^{-}_{4z}$\\
 & $D^{i}$ & $1$ & $-\mathrm{i}$ & $-1$ & $\mathrm{i}$\\
 & $(D^{i})^{-1}$ & $1$ & $\mathrm{i}$ & $-1$ & $-\mathrm{i}$\\
$r_{\alpha} = E$ & $D^{i}_{\alpha}$ & $1$ & $-\mathrm{i}$ & $-1$ & $\mathrm{i}$\\
$r_{\alpha} = C_{2x}$ & $D^{i}_{\alpha}$ & $1$ & $\mathrm{i}$ & $-1$ & $-\mathrm{i}$
\end{tabular}
\end{center}
条件只对 $r_{\alpha} = C_{2x}$ 满足。因此式 (4.6.10) 中的和成为 $(q_{i}/h)\chi^{{}^{1}E}(C^{2}_{2x}) = (2/2)\chi^{{}^{1}E}(E) = +1$。所以 $422\ (D_{4})$ 的表示 $E$ 是第一种——与我们在 2.4 节用直接方法得到的结论一致。

一个类似的例子是 4.5 节讨论过的四元数群（见式 (4.5.8)--(4.5.15)）。对 $\mathbf{Q}$，$\Gamma = {}^{1}E \uparrow \mathbf{Q}$ 的星有两个成员 ${}^{1}E$ 与 ${}^{2}E$，$h = 2$，而唯一满足条件 $D^{i}_{\alpha} = (D^{i})^{-1}$ 的 $r_{\alpha}$ 是 $Q$。此时式 (4.6.10) 的和为 $(q_{i}/h)\chi^{{}^{1}E}(Q^{2}) = (2/2)\chi^{{}^{1}E}(P^{2})$，由于 $P^{2} = Q^{2}$ 且 $\chi^{{}^{1}E}(P^{2}) = -1$，该和为 $-1$。于是 $\Gamma$ 是第二种表示（这正是四元数群不同于 $422\ (D_{4})$ 的地方）。

把式 (4.6.10) 给出的判别表示实性的简化条件用于空间群。若 $\Gamma^{i}_{j}$ 是小群 $\mathbf{G}^{\mathbf{k}}$ 的小表示 $\Gamma^{\mathbf{k}}_{j}$，则式 (4.6.10) 可用于确定 $\Delta^{\mathbf{k}}_{j}$ 的实性。总结为：

\noindent\textbf{定理 4.6.2}\quad $\Delta^{\mathbf{k}}_{j} = \Gamma^{\mathbf{k}}_{j} \uparrow \mathbf{G}$ 属于第一、第二或第三种，按照
\begin{equation*}
\frac{q_{\mathbf{k}}}{h}\sum_{R_{\alpha}}\chi^{\mathbf{k}}_{j}\left(\{R_{\alpha} \mid \mathbf{v}_{\alpha}\}^{2}\right) = +1,\ -1,\ \text{或}\ 0
\tag{4.6.11}
\end{equation*}
而定，其中 $\chi^{\mathbf{k}}_{j}$ 是 $\mathbf{k}$ 在 $\mathbf{G}$ 中小群的小表示 $j$ 的特征标，且求和限制在那些使 $\mathbf{k}$ 变为与 $-\mathbf{k}$ 等价的矢量的、$\mathbf{G}$ 关于 $\mathbf{T}$ 的陪集代表 $\{R_{\alpha} \mid \mathbf{v}_{\alpha}\}$ 上。

可以预期，式 (4.6.11) 中的和不依赖于对固定 $\alpha$ 选取哪个可能的陪集代表 $\{R_{\alpha} \mid \mathbf{v}_{\alpha}\}$；因为若 $\mathbf{t}$ 是布拉维格子的平移且 $R_{\alpha}\mathbf{k} \equiv -\mathbf{k}$，则可以证明 $\chi^{\mathbf{k}}_{j}$ 对两种选取取相同的值（利用小表示中平移的标量矩阵性质）。

确定点群或空间群表示的实性，其物理重要性在于：它可以用来判断把时间反演对称操作 $\theta$ 加到点群或空间群的操作上之后，是否会产生能级的额外简并（Herring 1937\textit{a}；另见 5.4、7.5 与 7.6 节）。
